\documentclass[11pt]{article}
\usepackage[margin=1in]{geometry}
\usepackage{amsmath,amssymb,amsthm}
\usepackage{hyperref}

\newtheorem{proposition}{Proposition}
\newtheorem{lemma}{Lemma}
\theoremstyle{definition}
\newtheorem{remark}{Remark}

\title{A Disproof of Conjecture \texttt{00000001238}\\
\large Log-concavity of Jacobian invariant factors, and the equality class}
\author{TLMC Submission}
\date{October 2026}

\begin{document}
\maketitle

\begin{abstract}
Conjecture \texttt{00000001238} asserts that (i) the invariant-factor sequence of the
Jacobian (sandpile group) of every graph is log-concave, and (ii) equality in the
concavity holds \emph{exactly} for direct sums of cycles.  Both clauses are false.
The complete bipartite graph $K_{2,4}$ has Jacobian
$\mathbb{Z}/2 \times \mathbb{Z}/2 \times \mathbb{Z}/8$, whose invariant-factor
sequence $(2,2,8)$ violates the middle concavity inequality ($2\cdot 2 = 4 < 2\cdot
8 = 16$); and $K_4$ has Jacobian $\mathbb{Z}/4 \times \mathbb{Z}/4$, whose sequence
$(4,4)$ \emph{has} equality although $K_4$ is not a direct sum of cycles (the edge
$\{0,1\}$ lies on the two triangles $\{0,1,2\}$ and $\{0,1,3\}$).  Moreover the two
clauses are jointly unsatisfiable even on the class the conjecture nominates: the
cactus $C_3 \vee C_3 \vee C_6$ is a direct sum of cycles yet has invariant factors
$(3,3,6)$, which are not log-concave ($9 < 18$).  An exhaustive sweep over all
$2^{15}$ graphs on six labeled vertices ($26{,}704$ connected) corroborates the
picture.  All numeric facts are certified in core Lean~4 with zero axioms and no
\texttt{sorry}, and reproduced by a standalone script.
\end{abstract}

\section{The conjecture}

\begin{quote}
``Definition: The invariant-factor sequence of the Jacobi group encodes the torsion
structure of the graph Laplacian. Conjecture: Log-concavity of the invariant-factor
sequence of the torsion group holds for all graphs; and equality in the concavity
holds exactly for direct sums of cycles. (Jacobian invariant factors log-concave)''
\end{quote}

\section{Conventions and reading}\label{sec:conv}

For a finite connected graph $G$ with vertices $V$, the \emph{Jacobian}
$\operatorname{Jac}(G)$ is the cokernel $\mathbb{Z}^{V}/\,\mathrm{rows}(L)$ of the
graph Laplacian $L$; equivalently $\mathbb{Z}^{|V|-1}/\mathrm{rows}(L^{\circ})$ for
any reduced Laplacian $L^{\circ}$ (delete one row and the matching column).  Its
\emph{invariant-factor sequence} is the nontrivial part $d_1 \mid d_2 \mid \dotsb \mid
d_k$ ($d_i > 1$) of the Smith normal form diagonal of $L^{\circ}$; this is the reading
fixed by the pipeline's own confirmed instance (SNF $(1,1,2,2,8)$, nontrivial factors
$[2,2,8]$, $2^2 = 4 < 2\cdot 8 = 16$).  \emph{Log-concavity} is the system of interior
inequalities $d_i^2 \ge d_{i-1} d_{i+1}$ for $1 < i < k$; \emph{equality} means every
interior inequality is an equality.  A \emph{direct sum of cycles} is a graph each of
whose blocks is a single edge or a cycle (a cactus: cycles joined at cut vertices),
whose Jacobian is then the direct sum of the $\mathbb{Z}/(\text{cycle length})$.

\begin{remark}[Convention robustness]\label{rem:robust}
Our counterexamples do not depend on these conventions.  The failure of $K_{2,4}$
occurs at an \emph{interior} triple $(2,2,8)$ with a strict inequality, so it
persists under the full-sequence reading (with leading $1$s), under endpoint
augmentation, and under reversal.  The failure of the ``exactly'' clause is witnessed
by $K_4$ (sequence $(4,4)$, equality $4\cdot 4 = 4 \cdot 4$) and, with interior
equalities, by $K_5$ (sequence $(5,5,5)$, all equalities $25 = 25$, $|\operatorname{Jac}(K_5)|
= 5^3 = 125$); neither is a cactus under any reading.  If instead ``direct sum of
cycles'' were read group-theoretically, the clause would be vacuous (every finite
abelian group is a direct sum of cyclic groups), so the graph-theoretic reading is
the only substantive one.
\end{remark}

\section{Counterexample A: log-concavity fails at $K_{2,4}$}

Take parts $A = \{0,1\}$, $B = \{2,3,4,5\}$.  Deleting vertex $0$ leaves the reduced
Laplacian
\[
L^{\circ}(K_{2,4}) \;=\; L_{24} \;=\;
\begin{pmatrix}
 4 & -1 & -1 & -1 & -1\\
-1 &  2 &  0 &  0 &  0\\
-1 &  0 &  2 &  0 &  0\\
-1 &  0 &  0 &  2 &  0\\
-1 &  0 &  0 &  0 &  2
\end{pmatrix}.
\]
The following integer matrices are unimodular, $\det U_{24} = \det V_{24} = -1$
(each is a product of elementary row, resp.\ column, operations applied to
$L_{24}$; all entries are verified by the kernel):
\[
U_{24} = \begin{pmatrix} 1&0&0&0&0\\ 0&0&1&0&0\\ 2&1&7&0&0\\ 2&1&6&1&0\\ 2&1&5&1&1 \end{pmatrix},
\qquad
V_{24} = \begin{pmatrix} 0&-1&0&0&2\\ -1&-4&1&0&1\\ 0&0&0&0&1\\ 0&0&-1&1&1\\ 0&0&0&-1&5 \end{pmatrix},
\]
\[
U_{24}\, L_{24}\, V_{24} \;=\; \operatorname{diag}(1,1,2,2,8).
\]
Hence
\[
\operatorname{Jac}(K_{2,4}) \;\cong\; \mathbb{Z}/1 \times \mathbb{Z}/1 \times
\mathbb{Z}/2 \times \mathbb{Z}/2 \times \mathbb{Z}/8 \;\cong\;
\mathbb{Z}/2 \times \mathbb{Z}/2 \times \mathbb{Z}/8 ,
\]
with invariant-factor sequence $(2,2,8)$.  As a sanity check, the order equals the
number of spanning trees,
$|\operatorname{Jac}(K_{2,4})| = 1\cdot 1\cdot 2\cdot 2\cdot 8 = 32 = 2^{4-1}\cdot
4^{2-1}$, which is the matrix-tree value $\tau(K_{m,n}) = m^{n-1} n^{m-1}$ at
$(m,n) = (2,4)$.  The middle concavity condition for $(2,2,8)$ requires
$d_2^{\,2} \ge d_1 d_3$, i.e.
\[
2 \cdot 2 = 4 \;\not\ge\; 2 \cdot 8 = 16 ,
\]
which fails.

\begin{proposition}\label{prop:A}
The invariant-factor sequence of $\operatorname{Jac}(K_{2,4})$ is $(2,2,8)$, which is
not log-concave.  Clause (i) of Conjecture \texttt{00000001238} is false.
\end{proposition}

\section{Counterexample B: equality without being a cycle sum}

\subsection{$K_4$}

The reduced Laplacian of $K_4$ (vertices $\{0,1,2,3\}$, vertex $0$ deleted) is
$L^{\circ}(K_4) = \mathrm{diag\text{-}off}\,(3,-1,-1;-1,3,-1;-1,-1,3)$, and
\[
U_4 = \begin{pmatrix} 1&0&0\\ 3&1&0\\ 2&1&1 \end{pmatrix},
\qquad
V_4 = \begin{pmatrix} 0&0&1\\ -1&1&1\\ 0&-1&2 \end{pmatrix},
\qquad
U_4\, L^{\circ}(K_4)\, V_4 = \operatorname{diag}(1,4,4),
\]
with $\det U_4 = \det V_4 = 1$.  So $\operatorname{Jac}(K_4) \cong
\mathbb{Z}/4 \times \mathbb{Z}/4$ with sequence $(4,4)$: the (length-two)
concavity condition $d_1^{\,2} \ge d_1 d_2$ reads $16 \ge 16$ and holds \emph{with
equality}.  But $K_4$ is not a direct sum of cycles:
\begin{itemize}
\item the edge $\{0,1\}$ lies on the two distinct triangles $\{0,1,2\}$ and
      $\{0,1,3\}$, while every edge of a direct sum of cycles lies on at most one
      cycle;
\item equivalently, an exhaustive check of all $2^6$ graphs on four vertices shows
      every connected cactus has at most $4$ edges and at most
      $\tau \in \{1,3,4\}$ spanning trees (there are exactly $31$ of them),
      whereas $K_4$ has $6$ edges and $\tau(K_4) = 4^{4-2} = 16$.
\end{itemize}
So the ``only if'' direction of clause (ii) fails.  The same failure in a form with
interior equalities is given by $K_5$:
$\operatorname{diag}(1,5,5,5)$, sequence $(5,5,5)$, all interior equalities
$25 = 25$, $|\operatorname{Jac}(K_5)| = 125 = 5^3$, and $K_5$ is very far from a
cactus.  (On six vertices, $K_6$ itself has sequence $(6,6,6,6)$ with all interior
equalities $36 = 36$.)

\begin{proposition}\label{prop:B}
$\operatorname{Jac}(K_4) \cong \mathbb{Z}/4 \times \mathbb{Z}/4$ has equality in the
concavity, but $K_4$ is not a direct sum of cycles.  Clause (ii) of Conjecture
\texttt{00000001238} is false.
\end{proposition}

\section{The two clauses are jointly unsatisfiable}

The conjecture's own equality class does not even satisfy clause (i).  The cactus
$C_3 \vee C_3 \vee C_6$ (three cycles joined at one vertex; a direct sum of cycles
in the conjecture's sense) has Jacobian
$\mathbb{Z}/3 \times \mathbb{Z}/3 \times \mathbb{Z}/6$ with sequence $(3,3,6)$, and
\[
3 \cdot 3 = 9 \;\not\ge\; 3 \cdot 6 = 18 .
\]
Thus no interpretation makes both clauses true: any graph class witnessing
``exactly'' must contain cacti, and cacti already break log-concavity.  (A caution
for the reading: the wedge $C_3 \vee C_4$ has Jacobian $\mathbb{Z}/3 \times
\mathbb{Z}/4 \cong \mathbb{Z}/12$ by the Chinese remainder theorem, since
$\gcd(3,4) = 1$; its invariant-factor sequence is the \emph{single} factor $(12)$.
Divisibility of the cycle lengths controls when a cactus Jacobian stays decomposed.)

\section{Exhaustive corroboration on six vertices}

Sweeping all $2^{15}$ graphs on six labeled vertices and computing the Smith normal
form of each reduced Laplacian by determinantal divisors (gcds of all $k \times k$
minors):
\begin{center}
\begin{tabular}{lr}
\hline
connected graphs on $6$ labeled vertices & $26{,}704$\\
nontrivial sequences of length $\ge 2$ that violate an interior inequality & $75$\\
graphs with torsion exactly $\mathbb{Z}/2 \times \mathbb{Z}/2 \times \mathbb{Z}/8$ & $15$
  \quad ($K_{2,4}$ among them)\\
equality sequences of length $\ge 2$ on non-cacti & $3{,}793$\\
\quad of which of length $\ge 3$ (e.g.\ $K_6$, $(6,6,6,6)$) & $31$\\
\hline
\end{tabular}
\end{center}
So both failure modes are ubiquitous already on six vertices.

\section{Machine verification}

The directory \texttt{lean4/} contains a core-only Lean~4 project (toolchain
\texttt{leanprover/lean4:v4.33.1}) certifying, with \textbf{zero axioms and no
\texttt{sorry}} (\texttt{Check.lean} audits all eleven theorems; none depends on any
axiom):
\begin{itemize}
\item \texttt{Tlmc1238.snf\_L24}: $U_{24} L_{24} V_{24} = \operatorname{diag}(1,1,2,2,8)$,
      $\det U_{24} = \det V_{24} = -1$;
\item \texttt{Tlmc1238.snf\_LK4}: $U_4 L^{\circ}(K_4) V_4 = \operatorname{diag}(1,4,4)$,
      $\det U_4 = \det V_4 = 1$;
\item \texttt{Tlmc1238.lap24\_is\_minor}, \texttt{lapK4\_is\_minor}: the two matrices
      equal the reduced Laplacians computed from the edge bitmasks $510$ and $615$;
\item \texttt{Tlmc1238.K4\_edge\_in\_two\_triangles}: both triangles through the edge
      $\{0,1\}$ exist and have distinct edge sets;
\item \texttt{Tlmc1238.seq\_228\_not\_logconcave}, \texttt{seq\_44\_logconcave\_eq},
      \texttt{seq\_555\_equalities}, \texttt{seq\_336\_not\_logconcave},
      \texttt{violation\_arith};
\item \texttt{Tlmc1238.refute}: the conjunction of all failures of both clauses.
\end{itemize}
Because core lemmas in this toolchain pull \texttt{propext} through
\texttt{Nat.testBit}, \texttt{Nat.lor} and overlapping match patterns, the
development uses only kernel-reduced arithmetic (\texttt{rfl}/\texttt{decide}),
structural \texttt{div}/\texttt{mod} bit tests, and disjoint pattern matches.
The script \texttt{reproduce.py} (pure standard library, ${\sim}15$\,s) recomputes
every number here --- unimodular products, determinantal divisors, the
four-vertex cactus classification, and the full six-vertex sweep.

\section{Scope of the disproof}

\begin{remark}
We refute the literal two-clause statement of Conjecture \texttt{00000001238} at the
concrete graphs $K_{2,4}$, $K_4$ (with $K_5$ and $C_3 \vee C_3 \vee C_6$ as
reinforcements).  The Smith theory of Laplacians, the matrix-tree theorem, and the
cactus decomposition of Jacobians of cycle sums are of course true and are used as
tools; the log-concavity claim and its ``exactly cycles'' equality claim are false
as stated.
\end{remark}

\end{document}
